\section{命令注册表与分发}

\subsection{注册表：静态指针数组}

\codefile{src/kernel/core/shell.c}
\begin{lstlisting}[style=cstyle]
// shell.c 中
extern const cmd_t cmd_cls;
extern const cmd_t cmd_shutdown;
extern const cmd_t cmd_cmds;
// ... 每加一个命令就加一行 extern ...

static const cmd_t* g_cmds[] = {
    &cmd_cls,
    &cmd_shutdown,
    &cmd_cmds,
    &cmd_mmap,
    &cmd_free,
    &cmd_pmm,
    &cmd_vmm,
    &cmd_heap,
    &cmd_vma,
};
\end{lstlisting}

\begin{figure}[H]
  \centering
  % figures/ch06/fig03-g_cmds-cmd_t.tex — auto-extracted from chapter source
\begin{tikzpicture}[node distance=0cm]
  % g_cmds array
  \foreach \i/\name in {0/cmd\_cls, 1/cmd\_shutdown, 2/cmd\_cmds, 3/cmd\_mmap, 4/{\ldots}} {
    \node[memcell, minimum width=3cm, minimum height=0.5cm] (g\i) at (0, -\i*0.55) {};
    \node[font=\tiny\ttfamily] at (g\i) {\name};
    \node[memaddr, left=0.1cm of g\i, font=\tiny\ttfamily] {[\i]};
  }
  \node[font=\small\ttfamily, above=0.2cm of g0] {g\_cmds[]};

  % cmd_t structures on the right
  \node[draw, minimum width=2.5cm, fill=blue!10, right=2.5cm of g0, font=\tiny\ttfamily] (cls) {name="cls", fn={\ldots}};
  \node[draw, minimum width=2.5cm, fill=orange!10, right=2.5cm of g1, font=\tiny\ttfamily] (shut) {name="shutdown", fn={\ldots}};

  \draw[pointer] (g0.east) -- (cls.west);
  \draw[pointer] (g1.east) -- (shut.west);
\end{tikzpicture}

  \caption{\texttt{g\_cmds[]} 是指向各 \texttt{cmd\_t} 常量的指针数组}
\end{figure}

\subsection{分发逻辑}

\codefile{src/kernel/core/shell.c}
\begin{lstlisting}[style=cstyle]
static void shell_dispatch(char* line) {
    if (!line || !*line) return;

    char* argv[16];
    int argc = shell_tokenize(line, argv, 16);
    if (argc <= 0) return;

    // 遍历注册表，匹配 name
    for (size_t i = 0; i < sizeof(g_cmds)/sizeof(g_cmds[0]); i++) {
        if (streq(argv[0], g_cmds[i]->name)) {
            g_cmds[i]->fn(argc, argv);
            return;
        }
    }

    printk("Unknown command: %s\n", argv[0]);
}
\end{lstlisting}

\begin{thinkbox}
当前遍历是 $O(n)$——对几十个命令完全够用。如果将来命令数量上千
（不太可能在内核态），可以改为排序数组 + 二分查找或哈希表。
但这是过度工程化——\term{YAGNI}（You Aren't Gonna Need It）。
\end{thinkbox}

% ============================================================================

